% GENERATED FILE. Edit canonical lessons, then run npm run generate:books.
% canonical-source-hash: fnv1a64:fe83da5bc41c28d3
% canonical-lessons: LA-C466-annona, LA-C466-caritas, LA-C466-vilitas, LA-C466-aestimo, LA-C466-veneo

\chapter{Corn supply, Dearness, Cheapness, To value, To be sold}
\label{ch:la-corn-supply-dearness-cheapness-to-value-to-be-sold}
\hlchaptermodality{466}

\begin{chapteropening}
\textbf{By the end of this chapter:} \emph{I can talk about prices: the corn supply, high prices and low prices, valuing goods and things being sold.}

Complete the last lesson of chapter 466: I can talk about prices: the corn supply, high prices and low prices, valuing goods and things being sold.
\end{chapteropening}

\section[annōna, annōnae]{annōna, annōnae — the corn supply; the price of grain}

\label{lesson:LA-C466-annona}



\begin{quote}
Before the new one: say the Latin for a quarter-as, then the Latin for a money-changer.
\end{quote}

\subsection*{You'll want to know: annōna, annōnae}
\textbf{annōna, annōnae} — \textquotedblleft{}the corn supply; the price of grain\textquotedblright{}.

\begin{quote}
Hōc annō annōna cāra est, quod pluvia nūlla fuit.
\end{quote}

\begin{quote}
\emph{This year grain is dear, because there was no rain.}
\end{quote}

First declension, feminine: \emph{annōna, annōnae f.} First the year's harvest, then the city's grain supply and its price.

\begin{cousinweb}
\emph{annōna} is built on \emph{annus} 'year' (EtymDB; Lewis \& Short): the produce of the year, and so the grain supply and its price.
\end{cousinweb}

\begin{grammarlens}[title={What you've built}]
One more word for telling, explaining and describing at B1.
\end{grammarlens}

\subsection*{Your turn}
\begin{itemize}
\raggedright
  \item \emph{Say it:} \emph{annōna, annōnae}
  \item \emph{Say it:} \emph{annōna, annōnae}, once more
  \item \emph{Say it:} a sentence of your own with \emph{annōna, annōnae}
  \item \emph{Recall:} say the Latin for a quarter-as, then the Latin for a money-changer, then say \emph{annōna, annōnae} again
\end{itemize}

\begin{tcolorbox}[breakable,colback=teal!4,colframe=teal!35!black,title={Before you move on}]
What does annōna, annōnae mean? (\textquotedblleft{}The corn supply; the price of grain\textquotedblright{}.) Say it once more.
\end{tcolorbox}
\section[cāritās, cāritātis]{cāritās, cāritātis — dearness, high prices; affection}

\label{lesson:LA-C466-caritas}



\begin{quote}
Before the new one: say the Latin for a money-changer, then the Latin for the corn supply.
\end{quote}

\subsection*{You'll want to know: cāritās, cāritātis}
\textbf{cāritās, cāritātis} — \textquotedblleft{}dearness, high prices; affection\textquotedblright{}.

\begin{quote}
Propter cāritātem annōnae pauperēs parum cibī habent.
\end{quote}

\begin{quote}
\emph{Because of the high price of grain the poor have too little food.}
\end{quote}

Third declension, feminine: \emph{cāritās, cāritātis f.}, from \emph{cārus} 'dear'. Dear in price, and dear to the heart. English \emph{charity} comes from it.

\begin{cousinweb}
\emph{cāritās} is built on \emph{cārus} 'dear, costly', from PIE \emph{keh\textsubscript{2}-} 'to desire' (EtymDB; Lewis \& Short). English \emph{charity} comes from it, in its sense 'affection'.
\end{cousinweb}

\begin{grammarlens}[title={What you've built}]
One more word for telling, explaining and describing at B1.
\end{grammarlens}

\subsection*{Your turn}
\begin{itemize}
\raggedright
  \item \emph{Say it:} \emph{cāritās, cāritātis}
  \item \emph{Say it:} \emph{cāritās, cāritātis}, once more
  \item \emph{Say it:} a sentence of your own with \emph{cāritās, cāritātis}
  \item \emph{Recall:} say the Latin for a money-changer, then the Latin for the corn supply, then say \emph{cāritās, cāritātis} again
\end{itemize}

\begin{tcolorbox}[breakable,colback=teal!4,colframe=teal!35!black,title={Before you move on}]
What does cāritās, cāritātis mean? (\textquotedblleft{}Dearness, high prices; affection\textquotedblright{}.) Say it once more.
\end{tcolorbox}
\section[vīlitās, vīlitātis]{vīlitās, vīlitātis — cheapness, low prices}

\label{lesson:LA-C466-vilitas}



\begin{quote}
Before the new one: say the Latin for the corn supply, then the Latin for dearness.
\end{quote}

\subsection*{You'll want to know: vīlitās, vīlitātis}
\textbf{vīlitās, vīlitātis} — \textquotedblleft{}cheapness, low prices\textquotedblright{}.

\begin{quote}
Post bonum annum vīlitās est, et agricolae minus pecūniae habent.
\end{quote}

\begin{quote}
\emph{After a good year prices are low, and the farmers have less money.}
\end{quote}

Third declension, feminine: \emph{vīlitās, vīlitātis f.} The opposite of \emph{cāritās}. Also worthlessness.

\begin{cousinweb}
\emph{vīlitās} is built on \emph{vīlis} 'cheap, of little value' with the suffix \emph{-itās} (EtymDB; Lewis \& Short). English \emph{vile} comes from \emph{vīlis}.
\end{cousinweb}

\begin{grammarlens}[title={What you've built}]
One more word for telling, explaining and describing at B1.
\end{grammarlens}

\subsection*{Your turn}
\begin{itemize}
\raggedright
  \item \emph{Say it:} \emph{vīlitās, vīlitātis}
  \item \emph{Say it:} \emph{vīlitās, vīlitātis}, once more
  \item \emph{Say it:} a sentence of your own with \emph{vīlitās, vīlitātis}
  \item \emph{Recall:} say the Latin for the corn supply, then the Latin for dearness, then say \emph{vīlitās, vīlitātis} again
\end{itemize}

\begin{tcolorbox}[breakable,colback=teal!4,colframe=teal!35!black,title={Before you move on}]
What does vīlitās, vīlitātis mean? (\textquotedblleft{}Cheapness, low prices\textquotedblright{}.) Say it once more.
\end{tcolorbox}
\section[aestimō, aestimāre]{aestimō, aestimāre — to value, to put a price on; to judge}

\label{lesson:LA-C466-aestimo}



\begin{quote}
Before the new one: say the Latin for dearness, then the Latin for cheapness.
\end{quote}

\subsection*{You'll want to know: aestimō, aestimāre}
\textbf{aestimō, aestimāre} — \textquotedblleft{}to value, to put a price on; to judge\textquotedblright{}.

\begin{quote}
Mercātor equum centum dēnāriīs aestimat.
\end{quote}

\begin{quote}
\emph{The merchant values the horse at a hundred denarii.}
\end{quote}

\emph{aestimō, aestimāre, aestimāvī, aestimātum}, with the price in the ablative or genitive: \emph{magnī aestimāre}, to value highly. English \emph{estimate} and \emph{esteem} come from it.

\begin{cousinweb}
\emph{aestimō} comes from older \emph{aestumāre} (EtymDB), and Lewis \& Short link it with \emph{aes} 'bronze, money', a valuing in money; the second part is not settled. English \emph{estimate} and \emph{esteem} come from it.
\end{cousinweb}

\begin{grammarlens}[title={What you've built}]
One more word for telling, explaining and describing at B1.
\end{grammarlens}

\subsection*{Your turn}
\begin{itemize}
\raggedright
  \item \emph{Say it:} \emph{aestimō, aestimāre}
  \item \emph{Say it:} \emph{aestimō, aestimāre}, once more
  \item \emph{Say it:} a sentence of your own with \emph{aestimō, aestimāre}
  \item \emph{Recall:} say the Latin for dearness, then the Latin for cheapness, then say \emph{aestimō, aestimāre} again
\end{itemize}

\begin{tcolorbox}[breakable,colback=teal!4,colframe=teal!35!black,title={Before you move on}]
What does aestimō, aestimāre mean? (\textquotedblleft{}To value, to put a price on; to judge\textquotedblright{}.) Say it once more.
\end{tcolorbox}
\section[vēneō, vēnīre]{vēneō, vēnīre — to be sold, to go on sale}

\label{lesson:LA-C466-veneo}



\begin{quote}
Before the new one: say the Latin for cheapness, then the Latin for to value.
\end{quote}

\subsection*{You'll want to know: vēneō, vēnīre}
\textbf{vēneō, vēnīre} — \textquotedblleft{}to be sold, to go on sale\textquotedblright{}.

\begin{quote}
In forō servī et equī vēneunt.
\end{quote}

\begin{quote}
\emph{Slaves and horses are sold in the forum.}
\end{quote}

Irregular like \emph{eō}: \emph{vēneō, vēnīre, vēniī}. It serves as the passive of \emph{vendō}. Do not confuse \emph{vēnīre} with \emph{venīre} 'to come'.

\begin{cousinweb}
\emph{vēneō} comes from \emph{vēnum eō} 'I go to sale' (EtymDB; Lewis \& Short), with \emph{vēnum} 'sale' and \emph{eō} 'to go'; \emph{vendō} 'to sell' is from \emph{vēnum dō} 'I give to sale'.
\end{cousinweb}

\begin{grammarlens}[title={What you've built}]
One more word for telling, explaining and describing at B1.
\end{grammarlens}

\subsection*{Your turn}
\begin{itemize}
\raggedright
  \item \emph{Say it:} \emph{vēneō, vēnīre}
  \item \emph{Say it:} \emph{vēneō, vēnīre}, once more
  \item \emph{Say it:} a sentence of your own with \emph{vēneō, vēnīre}
  \item \emph{Recall:} say the Latin for cheapness, then the Latin for to value, then say \emph{vēneō, vēnīre} again
\end{itemize}

\begin{tcolorbox}[breakable,colback=teal!4,colframe=teal!35!black,title={Before you move on}]
What does vēneō, vēnīre mean? (\textquotedblleft{}To be sold, to go on sale\textquotedblright{}.) Say it once more.
\end{tcolorbox}
